\documentclass[a4paper]{article}

\usepackage[T1]{fontenc}
\usepackage[italian]{babel}
\usepackage{tikz}
\usepackage{pgfplots}
\usepackage{amsmath}

\usetikzlibrary{calc}



\begin{document}

\title{Travi reticolari}
\author{Andrea Marchi}

\maketitle

\textcolor{red}{Per ogni trave fare confronto in termini di:
\begin{itemize}
\item abbassamento in mezzeria
\item massimo sforzo normale nelle aste
\item istogramma con lo sforzo normale nelle aste (ordinato)
\item rapporto tra aste compresse e aste tese
\end{itemize}
}






\section{Trave francese (Howe)}

Questo tipo di trave reticolare e` caratterizzato dall'avere le diagonali tese. Il vantaggio e` che le aste di parete compresse sono quelle verticali, le quali hanno una lunghezza minore rispetto a quelle diagonali e, quindi, esibiscono instabilita` piu` difficilmente.

\begin{tikzpicture}[scale=0.6]
% Parametri visualizzazione:
\def\dispscale{2000}
% Nodi indeformati:
\node[draw, circle, fill=white, inner sep=2pt, label={[above right]0}] (N0) at (0.0,0) {};
\node[draw, circle, fill=white, inner sep=2pt, label={[above right]1}] (N1) at (2.0,0) {};
\node[draw, circle, fill=white, inner sep=2pt, label={[above right]2}] (N2) at (4.0,0) {};
\node[draw, circle, fill=white, inner sep=2pt, label={[above right]3}] (N3) at (6.0,0) {};
\node[draw, circle, fill=white, inner sep=2pt, label={[above right]4}] (N4) at (8.0,0) {};
\node[draw, circle, fill=white, inner sep=2pt, label={[above right]5}] (N5) at (10.0,0) {};
\node[draw, circle, fill=white, inner sep=2pt, label={[above right]6}] (N6) at (12.0,0) {};
\node[draw, circle, fill=white, inner sep=2pt, label={[above right]7}] (N7) at (14.0,0) {};
\node[draw, circle, fill=white, inner sep=2pt, label={[above right]8}] (N8) at (16.0,0) {};
\node[draw, circle, fill=white, inner sep=2pt, label={[above right]9}] (N9) at (2.0,-2.0) {};
\node[draw, circle, fill=white, inner sep=2pt, label={[above right]10}] (N10) at (4.0,-2.0) {};
\node[draw, circle, fill=white, inner sep=2pt, label={[above right]11}] (N11) at (6.0,-2.0) {};
\node[draw, circle, fill=white, inner sep=2pt, label={[above right]12}] (N12) at (8.0,-2.0) {};
\node[draw, circle, fill=white, inner sep=2pt, label={[above right]13}] (N13) at (10.0,-2.0) {};
\node[draw, circle, fill=white, inner sep=2pt, label={[above right]14}] (N14) at (12.0,-2.0) {};
\node[draw, circle, fill=white, inner sep=2pt, label={[above right]15}] (N15) at (14.0,-2.0) {};
% Aste indeformate:
\draw[black, thin] (N0) -- (N1);
\draw[black, thin] (N1) -- (N2);
\draw[black, thin] (N2) -- (N3);
\draw[black, thin] (N3) -- (N4);
\draw[black, thin] (N4) -- (N5);
\draw[black, thin] (N5) -- (N6);
\draw[black, thin] (N6) -- (N7);
\draw[black, thin] (N7) -- (N8);
\draw[black, thin] (N9) -- (N10);
\draw[black, thin] (N10) -- (N11);
\draw[black, thin] (N11) -- (N12);
\draw[black, thin] (N12) -- (N13);
\draw[black, thin] (N13) -- (N14);
\draw[black, thin] (N14) -- (N15);
\draw[black, thin] (N1) -- (N9);
\draw[black, thin] (N2) -- (N10);
\draw[black, thin] (N3) -- (N11);
\draw[black, thin] (N4) -- (N12);
\draw[black, thin] (N5) -- (N13);
\draw[black, thin] (N6) -- (N14);
\draw[black, thin] (N7) -- (N15);
\draw[black, thin] (N0) -- (N9);
\draw[black, thin] (N1) -- (N10);
\draw[black, thin] (N2) -- (N11);
\draw[black, thin] (N3) -- (N12);
\draw[black, thin] (N5) -- (N12);
\draw[black, thin] (N6) -- (N13);
\draw[black, thin] (N7) -- (N14);
\draw[black, thin] (N8) -- (N15);
% Spostamenti nodi:
\node (u0) at (0.00000000000000000000,0.00000000000000000000) {};
\node (u1) at (-0.00000761904761904765,-0.00012059753999806604) {};
\node (u2) at (-0.00002068027210884360,-0.00021762299373818096) {};
\node (u3) at (-0.00003700680272108842,-0.00028019200747884818) {};
\node (u4) at (-0.00005442176870748289,-0.00030286240434932013) {};
\node (u5) at (-0.00007183673469387734,-0.00028019200747884812) {};
\node (u6) at (-0.00008816326530612215,-0.00021762299373818058) {};
\node (u7) at (-0.00010122448979591799,-0.00012059753999806560) {};
\node (u8) at (-0.00010884353741496554,0.00000000000000000000) {};
\node (u9) at (-0.00009142857142857122,-0.00011297849237901842) {};
\node (u10) at (-0.00008380952380952358,-0.00021218081686743279) {};
\node (u11) at (-0.00007074829931972766,-0.00027692670135639921) {};
\node (u12) at (-0.00005442176870748287,-0.00030068553360102076) {};
\node (u13) at (-0.00003809523809523802,-0.00027692670135639916) {};
\node (u14) at (-0.00002503401360544215,-0.00021218081686743222) {};
\node (u15) at (-0.00001741496598639455,-0.00011297849237901801) {};
\node[draw, circle, fill=black, inner sep=2pt] (D0) at ($ (N0) + \dispscale*(u0) $) {};
\node[draw, circle, fill=black, inner sep=2pt] (D1) at ($ (N1) + \dispscale*(u1) $) {};
\node[draw, circle, fill=black, inner sep=2pt] (D2) at ($ (N2) + \dispscale*(u2) $) {};
\node[draw, circle, fill=black, inner sep=2pt] (D3) at ($ (N3) + \dispscale*(u3) $) {};
\node[draw, circle, fill=black, inner sep=2pt] (D4) at ($ (N4) + \dispscale*(u4) $) {};
\node[draw, circle, fill=black, inner sep=2pt] (D5) at ($ (N5) + \dispscale*(u5) $) {};
\node[draw, circle, fill=black, inner sep=2pt] (D6) at ($ (N6) + \dispscale*(u6) $) {};
\node[draw, circle, fill=black, inner sep=2pt] (D7) at ($ (N7) + \dispscale*(u7) $) {};
\node[draw, circle, fill=black, inner sep=2pt] (D8) at ($ (N8) + \dispscale*(u8) $) {};
\node[draw, circle, fill=black, inner sep=2pt] (D9) at ($ (N9) + \dispscale*(u9) $) {};
\node[draw, circle, fill=black, inner sep=2pt] (D10) at ($ (N10) + \dispscale*(u10) $) {};
\node[draw, circle, fill=black, inner sep=2pt] (D11) at ($ (N11) + \dispscale*(u11) $) {};
\node[draw, circle, fill=black, inner sep=2pt] (D12) at ($ (N12) + \dispscale*(u12) $) {};
\node[draw, circle, fill=black, inner sep=2pt] (D13) at ($ (N13) + \dispscale*(u13) $) {};
\node[draw, circle, fill=black, inner sep=2pt] (D14) at ($ (N14) + \dispscale*(u14) $) {};
\node[draw, circle, fill=black, inner sep=2pt] (D15) at ($ (N15) + \dispscale*(u15) $) {};
% Aste deformate:
\draw[very thick] (D0) -- (D1);
\draw[very thick] (D1) -- (D2);
\draw[very thick] (D2) -- (D3);
\draw[very thick] (D3) -- (D4);
\draw[very thick] (D4) -- (D5);
\draw[very thick] (D5) -- (D6);
\draw[very thick] (D6) -- (D7);
\draw[very thick] (D7) -- (D8);
\draw[very thick] (D9) -- (D10);
\draw[very thick] (D10) -- (D11);
\draw[very thick] (D11) -- (D12);
\draw[very thick] (D12) -- (D13);
\draw[very thick] (D13) -- (D14);
\draw[very thick] (D14) -- (D15);
\draw[very thick] (D1) -- (D9);
\draw[very thick] (D2) -- (D10);
\draw[very thick] (D3) -- (D11);
\draw[very thick] (D4) -- (D12);
\draw[very thick] (D5) -- (D13);
\draw[very thick] (D6) -- (D14);
\draw[very thick] (D7) -- (D15);
\draw[very thick] (D0) -- (D9);
\draw[very thick] (D1) -- (D10);
\draw[very thick] (D2) -- (D11);
\draw[very thick] (D3) -- (D12);
\draw[very thick] (D5) -- (D12);
\draw[very thick] (D6) -- (D13);
\draw[very thick] (D7) -- (D14);
\draw[very thick] (D8) -- (D15);
\end{tikzpicture}





\section{Espressioni approssimate}

\subsection{Abbassamento in mezzeria}

Vorrei verificare se la freccia in mezzeria e` simile a quello di una trave semplicemente appoggiata con un carico concentrato in mezzeria pari al carico totale. Una trave per cui la sezione e` pari a due aree pari alle aree dei correnti distanti quanto l'altezza della trave.
Per una sezione simile l'inerzia e' pari a
$$ I = A \left(\frac{H}{2}\right)^2 + A \left(\frac{H}{2}\right)^2 = \frac{AH^2}{2} $$
L'abbassamento in mezzeria di una trave semplicemente appoggiata con carico concentrato $P$ in mezzeria e`
$$ f = \frac{PL^3}{48EI} = \frac{pL^4}{24EAH^2} $$
in cui abbiamo sostituito l'inerzia trovata in precedenza e il carico distribuito $p = P/L$.

\textcolor{red}{Mi viene in mente che questo approccio non tiene conto della deformabilita` delle aste di parete. Infatti dipende solamente dall'area dei correnti... Forse funziona solamente se le aste di parete hanno un'area particolare...}

per la seguente analisi i dati sono
\begin{itemize}
\item lunghezza $L=16$ m
\item peso distribuito $p=10$ kN/m
\item $E=210$ GPa
\item $A=0.1$ m$^2$
\end{itemize}
\begin{tikzpicture}
\begin{axis}[
    title={Abbassamento mezzeria trave Howe ($L=16$)},
    xlabel={Numero elementi corrente superiore},
    ylabel={$\frac{\text{Abbassamento}}{\frac{pL^4}{24EAH^2}}$ [-]},
%    xmin=0, xmax=100,
%    ymin=0, ymax=120,
    xtick={4,6,8,10,12,14,16,18,20},
%    ytick={0,20,40,60,80,100,120},
    legend pos=north west,
    ymajorgrids=true,
    grid style=dashed,
]

\addplot[
	color=black,
	mark=square,
	]
	coordinates {
    ( 4, 1.153682)
	( 6, 0.894344)
	( 8, 0.805197)
	(10, 0.763126)
	(12, 0.739849)
	(14, 0.725764)
	(16, 0.716800)
	(18, 0.710963)
	(20, 0.707168)
	};

\addplot[
	color=black,
	mark=o,
	]
	coordinates {
    ( 4, 1.26344812)
	( 6, 1.00885416)
	( 8, 0.93165681)
	(10, 0.90444171)
	(12, 0.89760085)
	(14, 0.90093923)
	(16, 0.91006674)
	(18, 0.92279501)
	(20, 0.93791425)
	};

\addplot[
	color=black,
	mark=diamond,
	]
	coordinates {
    ( 4, 1.49511718)
	( 6, 1.26589151)
	( 8, 1.22919419)
	(10, 1.24837239)
	(12, 1.29092893)
	(14, 1.34541756)
	(16, 1.40680814)
	(18, 1.47254932)
	(20, 1.54121004)
	};
  
\legend{$H = 1$ m, $H = 2$ m, $H = 3$ m}
	
\end{axis}
\end{tikzpicture}

La trave con $H=3$ m ha un rapporto lunghezza/altezza di 5.33.






\end{document}